\documentclass[10pt]{article}
\providecommand{\FIGDIR}{fig}
\newcommand{\DIGESTDATE}{31 October 2026}
\input{preamble}

\begin{document}

% ==================================================================== MASTHEAD
\thispagestyle{empty}
\begin{tikzpicture}[remember picture, overlay]
  \fill[Deep] (current page.north west) rectangle
      ($(current page.north east)+(0,-67mm)$);
  \fill[Coral] ($(current page.north west)+(0,-67mm)$) rectangle
      ($(current page.north east)+(0,-68.4mm)$);
  \foreach \x/\y/\r in {22/12/0.9, 41/26/1.5, 63/9/0.7, 88/31/1.1, 112/17/1.9,
                        138/29/0.8, 159/13/1.3, 178/54/1.0, 196/21/1.6,
                        30/44/0.6, 74/57/1.2, 125/48/0.9, 168/6/0.8, 191/59/0.7,
                        52/18/1.0, 100/60/0.6, 148/20/1.1, 66/49/1.4, 205/45/0.8}
    \fill[Sky, opacity=0.30] ($(current page.north west)+(\x mm,-\y mm)$) circle (\r mm);
\end{tikzpicture}

\vspace*{4mm}
{\sffamily\fontsize{7.6}{9}\selectfont\color{Coral}%
 MONTHLY ROLLUP\;\textbullet\;PhD RESEARCH WATCH\par}
\vspace{2.2mm}
{\sffamily\bfseries\fontsize{27}{29}\selectfont\color{white}%
 Particulate Matter\par}
{\sffamily\fontsize{27}{29}\selectfont\color{Sky}%
 Monitoring \& Health Effects\par}
\vspace{5.0mm}
{\sffamily\fontsize{8.4}{10}\selectfont\color{white}%
 October 2026\;\;\textcolor{Coral}{\textbar}\;\;Entry window 1--31 October 2026, 31 daily issues\;\;%
 \textcolor{Coral}{\textbar}\;\;523 records pooled\par}
\vspace{18mm}

% -------------------------------------------------------------- metric strip
\noindent
\begin{minipage}[t]{0.190\linewidth}\metric{523}{RECORDS POOLED\\(31 DAILY ISSUES)}{Deep}\end{minipage}\hfill
\begin{minipage}[t]{0.190\linewidth}\metric{130}{EXPOSURE ASSESSMENT\\(25\%), LARGEST BAND}{Teal}\end{minipage}\hfill
\begin{minipage}[t]{0.190\linewidth}\metric{113}{EFFECT ESTIMATES\\(113 WITH A CI)}{Sky}\end{minipage}\hfill
\begin{minipage}[t]{0.190\linewidth}\metric{75}{TIER-A\\RECORDS (14\%)}{Amber}\end{minipage}\hfill
\begin{minipage}[t]{0.190\linewidth}\metric{112}{METADATA-ONLY\\RECORDS (21\%)}{Clay}\end{minipage}

\vspace{6mm}

% ------------------------------------------------------------ signal of month
\section*{\sffamily\bfseries\color{Deep}Signal of the month}
\vspace{-1.5mm}{\color{Coral}\rule{\linewidth}{1.1pt}}\vspace{2.5mm}

\begin{keybox}
{\sffamily\bfseries\small\color{Deep}October pulled the exposure variable apart at both ends: the
epidemiology moved from PM$_{2.5}$ mass to its components, its number concentration and its source,
while the measurement literature showed that the references every sensor is calibrated against carry
biases of their own.}\\[1.4mm]
\footnotesize
\textbf{The health side.} In five weeks the effect was carried by something other than mass in every large
cohort that could separate it. Never-smoker lung cancer in the China Kadoorie Biobank: nitrate HR
\textbf{1.35} (1.23--1.48) per IQR against 1.13 for the mixture (3 Oct). Dementia mortality in 2.1 million
adults: HR \textbf{1.22} (1.17--1.27) per 10,000\,cm$^{-3}$ of ultrafine particle number (5 Oct). Incident
ADRD in 39.7 million Medicare beneficiaries: HR \textbf{1.029} (1.024--1.035) per 1\,$\mu$g\,m$^{-3}$ of
smoke PM$_{2.5}$ but \textbf{1.090} (1.077--1.103) per 1\,$\mu$g\,m$^{-3}$ of smoke EC+OC (W42). eGFR decline in
307,639 adults at HR \textbf{1.27--1.30} for sulfate, nitrate and ammonium; ADHD following elemental carbon
at HR \textbf{1.24} independent of PM$_{2.5}$ (W43). Prenatal wildfire PM$_{2.5}$ and autism in 893,212
Medicaid pairs: HR \textbf{1.30} (1.07--1.59) per 0.75\,$\mu$g\,m$^{-3}$, ADHD null (\textit{JAMA Pediatr},
26 Oct). \textbf{Every one of these component, number and smoke exposures is modelled} --- CTM,
satellite--ML, smoke-fraction regression or source apportionment --- not measured.

\vspace{1mm}
\textbf{The measurement side.} A three-level meta-analysis of 96 paired comparisons put professional
photometers at \textbf{36\,\% high} against gravimetry: lnR 0.31 (0.25--0.37), \textbf{I$^2$ 78.4\,\%},
prediction interval 0.11--0.63; site calibration was the strongest moderator (factory 43\,\% high, site
13\,\%) and no eligible study tested a sub-\$500 sensor (\textit{Ann Work Expo Health}, 23 Oct).
Gravimetry itself retains strongly bound water through standard conditioning, a \textbf{+3 to +18\,\%}
mass bias largest for PM$_1$ (17 Oct). The PurpleAir housing warms the PMS5003 detection zone and dries
the sample before it is counted, so RH corrections fitted on bare sensors assume the wrong humidity (27 Oct).
\textbf{The month's verdict: a network that reports corrected mass is converging on a quantity the
epidemiology is leaving behind, and is converging on it against a reference that is itself biased by
composition and water.}
\end{keybox}

\vspace{3mm}

\section*{\sffamily\bfseries\color{Deep}What changed in October}
\vspace{-1.5mm}{\color{Coral}\rule{\linewidth}{1.1pt}}\vspace{2.5mm}

\footnotesize
\begin{itemize}
\item \textbf{Volume fell 16\,\% and the mix moved toward modelling and the lung.} 523 records against 622
in September, on Sunday--Saturday week totals of 59 (1--3 Oct, partial), 87, 123, 119 and 135. Daily volume
ran from \textbf{1} (Sunday 11 Oct) to \textbf{30} (1 and 27 Oct). \textit{Exposure assessment \&
modelling} grew in every week (16, 23, 26, 31, \textbf{34}) and is the largest band at 130 (25\,\%);
\textit{Respiratory \& allergic} rose from \textbf{1 to 17} per week. \textit{Sensing} fell to 77 (15\,\%,
from 19\,\%), \textit{Cardiovascular \& metabolic} 46 $\rightarrow$ 28 and \textit{Occupational \& indoor}
72 $\rightarrow$ 50. The measurement share (sensing plus exposure) read 42, 44, 39, 39 and \textbf{37\,\%}
by week: the measurement side is now mostly models, not instruments.

\item \textbf{Calibration papers started reporting the hold-out that matters --- September's prediction
(2) met.} Mercado-Velazquez et al.\ (\textit{Sensors}, 27 Oct) co-located low-, mid-cost and reference
instruments for ten months at two London sites: PM$_{2.5}$ calibrations trained April--October held
through winter and transferred between sites at R$^2\geq$0.56, while NO$_2$ and CO did not (R$^2\leq$0.28);
linear T/RH regression matched ML. Sylvester et al.\ (\textit{ACS ES\&T Air}, 14 Oct): the national
Barkjohn correction beat local fits in Philadelphia (RMSE 1.64\,$\mu$g\,m$^{-3}$, R$^2$ 0.683), and R$^2$
fell to 0.567 after a year in the field. Alturki et al.\ (22 Oct) reached spatially blocked R$^2$
0.82--0.84 for a GEOS-CF emulator --- and a random forest on latitude and longitude alone scored
\textbf{0.91}, which is the leakage argument in one number. A conformal PM$_{2.5}$ map lost nominal
coverage (0.794) under temporal hold-out and its preregistered hypothesis failed (8 Oct). Fang et al.'s
66-study review (10 Oct) gives the field's baseline: median calibrated R$^2$ 0.92 in the laboratory
against \textbf{0.72} in the field, median RMSE 2.6 against \textbf{12.3}\,$\mu$g\,m$^{-3}$.

\item \textbf{The reference became the object of study.} Beyond the photometer meta-analysis and the
filter-water bias: a fluorescence--Raman--Mie lidar that gates its extinction-to-mass regression on
aerosol type cut surface RMSE by 38\,\% and bias by $\sim$83\,\% under date-separated validation
(\textit{AMT}, 16 Oct); four prototype electrostatic mass monitors at one power-plant site agreed only
within pairwise ratios of 0.33--2.83 for PM$_{2.5}$ (single-unit ICC 0.42 for PM$_{10}$), with no reference to
say which was right (20 Oct); and a chassis-dynamometer tire-wear campaign saw its reference CPC's response
double over the campaign (15 Oct).

\item \textbf{Number concentration acquired endpoints.} Beside the dementia cohort: kidney-disease ED
visits in New York State rose 0.74\,\% (0.28--1.20) per IQR of CTM-modelled UFP at lag 0--3 across
269,945 visits (29 Oct); source-resolved lung deposition ranged from 76\,\% (new particle formation)
to 27\,\% (solid fuel) at equal ambient number (8 Oct); 3-D printing emits almost entirely below the
optical cut of consumer sensors (25 Oct). Ultrafine and size-resolved number records rose to 23
(4.4\,\%, from 3\,\%).

\item \textbf{Wildfire smoke separated from ambient PM$_{2.5}$ as its own exposure} --- with its own
composition (the EC+OC ADRD contrast), its own lags (a deconvolved US mortality response peaking three
years after a smoke pulse, 1.7 million deaths 2011--2023) and its own endpoints (head-and-neck cancer
RR 1.41, 1.24--1.60; autism). The month also produced informative nulls: incident depression in 179,218
All of Us participants (preregistered, 30 Oct) and ADHD in the Medicaid cohort. On the measurement side,
32 Idaho AQS sites over 20 years show fire-season PM$_{2.5}$ rising 2.54\,$\mu$g\,m$^{-3}$ per decade while
the annual mean stayed flat (31 Oct).

\item \textbf{Scale and inventories.} Yu et al.\ (\textit{Lancet Planet Health}, 12 Oct): 346.5 million
admissions in 949 communities, RR \textbf{1.080} (1.074--1.086) per 10\,$\mu$g\,m$^{-3}$, with malignant
neoplasms (1.108) and mental disorders (1.100) above respiratory (1.067); heterogeneity described as
substantial but not reported with the pooled estimate. China's BC burden was put at $>$540,000 deaths in
2015, twice prior estimates (\textit{Nat Commun}, 8 Oct), and road tire-and-road-wear PM$_{2.5}$ at 425
kt in 2023, projected to reach 67\,\% of traffic PM$_{2.5}$ emissions by 2050 (\textit{Nat Commun}, 9 Oct).

\item \textbf{One network measured the fraction the epidemiology points to.} Manchanda et al.\
(\textit{Sci Adv}, 28 Oct) assimilated dense mobile and fixed black-carbon measurements (6.7 sites
km$^{-2}$) into hourly 150\,m emission maps for West Oakland: the posterior raised the BC inventory
$\sim$33\,\% and quadrupled the exposure share of neighbourhood and port sources (5 $\rightarrow$
20\,\%); $\sim$10 fixed monitors plus mobile maps recovered the result. It is not a low-cost network with
a season-long field calibration, so September's prediction (1) is \textbf{not met} --- but it is the
closest record yet.
\end{itemize}

\clearpage

% ==================================================================== FIGURES
\section*{\sffamily\bfseries\color{Deep}Corpus analytics --- October 2026}
\vspace{-1.5mm}{\color{Coral}\rule{\linewidth}{1.1pt}}\vspace{2.5mm}

\noindent\includegraphics[width=\linewidth]{\FIGDIR/f1_subtopics.png}
\figcap{Subtopic distribution across all 523 pooled records; all ten subtopics populated.
\textit{Exposure assessment \& modelling} 130 (25\,\%) and \textit{Sensing, forecasting \&
instrumentation} 77 (15\,\%) together hold 40\,\% (September 43\,\%, August 48\,\%). \textit{Burden, policy \&
mitigation} 64, \textit{Mechanistic toxicology} 58, \textit{Occupational \& indoor} 50, \textit{Respiratory
\& allergic} 40, \textit{Cardiovascular \& metabolic} and \textit{Other clinical} 28 each, \textit{Neuro /
mental health} and \textit{Reproductive \& developmental} 24 each. Single-label assignment by dominant
contribution.}

\vspace{3mm}

\noindent\includegraphics[width=\linewidth]{\FIGDIR/w1_subtopic_trend.png}
\figcap{Records per Sunday--Saturday week by subtopic across the 31 October issues (the first week is
partial): 59 for 1--3 Oct, then 87, 123, 119 and 135. \textit{Exposure} rose in every week (16 to 34) and
\textit{Respiratory} from 1 to 17; \textit{Sensing} peaked at 22 in 11--17 Oct. Daily volume ranged from 1
(Sunday 11 Oct) to 30 and tracks publisher deposit timing. Issues for 3--8 Oct and 23--29 Oct were built
late on recovery runs; each kept its own entry-date window.}

\vspace{3mm}
\noindent\includegraphics[width=\linewidth]{\FIGDIR/f2a_architecture.png}
\figcap{Study architecture. Metadata-only 112, measurement campaigns 95, modelling/inventory 84, cohorts 49,
cross-sectional 40, reviews 37, experimental/toxicology 36, chamber/laboratory 23, acute designs 21,
ecological 17 (drawn as two bars, ``Observational -- ecological'' 11 and ``Ecological'' 6: a label-map
defect, not two designs), trials or interventions 8, sensor co-location 1. \textbf{Metadata-only rose
again, 12\,\% $\rightarrow$ 16\,\% $\rightarrow$ 21\,\%} across August--October. Trials: 8 in 523, unchanged in
absolute terms from September.}

\clearpage

%% f2b_endpoint.png DELIBERATELY OMITTED (as Aug/Sep): ENDPOINT_CANON is still unbounded at month scale.
\begin{keybox}
{\sffamily\bfseries\footnotesize\color{Deep}Health endpoint distribution --- not drawn}\\[1.2mm]
\footnotesize
The endpoint figure is omitted for the third month running: the canonical endpoint map is still not
capped, so at month scale the chart is several dozen rows at an unreadable size. The subtopic weekly
series (above) carries the shift that matters --- respiratory endpoints rising through the month and
cardiometabolic endpoints falling. \textbf{Fix still owed}: cap the map at $\sim$20 categories.
\end{keybox}

\vspace{3mm}

\noindent\includegraphics[width=\linewidth]{\FIGDIR/f3_metric_geography.png}
\figcap{Left --- particle metric. PM$_{2.5}$ alone 216, unspeciated PM or emissions 138, PM$_{2.5}$ with
PM$_{10}$ 73, optical properties 26, \textbf{size-resolved number 21 and ultrafine 2} (4.4\,\% together,
against 3\,\% in September), composition/speciation 18, PM$_{10}$ alone 17, bioaerosol/coarse 7,
respirable dust 3, composite indices 2. Composition is under-counted here: component exposures inside
cohort records are filed under PM$_{2.5}$. Right --- geography. Global / multi-region 135 (inflated by
metadata-only records), \textbf{China 120 (23\,\%, from 27\,\%)}, Europe 72, North America 69, South Asia
32, East Asia ex-China 24, Middle East \& N.\ Africa 20, Southeast Asia 18, \textbf{Sub-Saharan Africa 14}
(22 in September), Latin America 13, Central Asia 5, Oceania 1.}

\clearpage

\noindent\includegraphics[width=\linewidth]{\FIGDIR/f4_heatmap.png}
\figcap{Subtopic $\times$ study-architecture cross-tabulation. Dense cells: exposure $\times$ modelling,
sensing $\times$ measurement campaign, and metadata-only across every row. The trial column remains nearly
empty in every health subtopic, and the sensing and exposure rows still contribute almost nothing to the
effect-estimate pool.}

\vspace{3mm}

\noindent\includegraphics[width=\linewidth]{\FIGDIR/f6_lifecourse.png}
\figcap{Life-course windows across the month (non-exclusive; only records with an identifiable human
window): in utero 30, infancy/childhood 36, adolescence 11, working age 54, older adults 29 (September
41, 35, 10, 77, 40). \textbf{Adolescence remains the thinnest window} for the fourth month running.}

\clearpage

\noindent\includegraphics[width=\linewidth,height=0.80\textheight,keepaspectratio]{\FIGDIR/f5_forest.png}
\figcap{Quantitative effect estimates extracted across October, panel 1 of 5. \textbf{This is not a
meta-analysis and no pooled summary is drawn.} Contrasts are not commensurable (per 1, per 10, per IQR,
per 10,000\,cm$^{-3}$, threshold, quartile, fuel-type, smoke-fraction contrasts) and the outcomes are
unrelated, so \textbf{between-study heterogeneity is not estimable and none is reported}. The one pooled
estimate in the set carries its own: the photometer--gravimetric lnR 0.31 at \textbf{I$^2$ 78.4\,\%}; the
six-country admissions RR (Yu et al.) was published without its heterogeneity statistic. All 113 estimates
carry an interval. Of 97 on ratio scales (HR 45, OR 30, RR 19, IRR 3), \textbf{82 exclude the null in the
harmful direction, 2 fall below it and 13 cross it}.}

\clearpage
\noindent\includegraphics[width=\linewidth,height=0.86\textheight,keepaspectratio]{\FIGDIR/f5_forest_2.png}
\figcap{Effect estimates, panel 2 of 5. Same caveats as panel 1.}
\clearpage
\noindent\includegraphics[width=\linewidth,height=0.86\textheight,keepaspectratio]{\FIGDIR/f5_forest_3.png}
\figcap{Effect estimates, panel 3 of 5. Same caveats as panel 1.}
\clearpage
\noindent\includegraphics[width=\linewidth,height=0.86\textheight,keepaspectratio]{\FIGDIR/f5_forest_4.png}
\figcap{Effect estimates, panel 4 of 5. Same caveats as panel 1.}
\clearpage
\noindent\includegraphics[width=\linewidth,height=0.80\textheight,keepaspectratio]{\FIGDIR/f5_forest_5.png}
\figcap{Effect estimates, panel 5 of 5. Neither ratio estimate below the null is a protective PM effect:
IVF live birth OR 0.969 (0.954--0.984) is harm on a decreasing scale, and ragweed sensitisation OR 0.69
(0.50--0.96) is a protective co-exposure (rural residence). Of the 16 non-ratio estimates (FEV$_1$, eGFR,
years, PAF, \%-change, HAZ), 14 exclude zero. The largest ratios --- OR 4.35 (2.41--7.83) for VSD after
post-conception PM$_{10}$ in a Delhi case-control set, and HR 2.55 (1.78--3.67) for cervical-cancer death in
a 3,758-patient Fudan cohort --- come from single-centre designs with wide intervals and, in the Delhi
case, a non-significant window contrast.}

\clearpage

\section*{\sffamily\bfseries\color{Deep}Emerging, fading, and recurring}
\vspace{-1.5mm}{\color{Coral}\rule{\linewidth}{1.1pt}}\vspace{2.5mm}

\footnotesize
\textbf{Emerging.} \textbf{(i) The exposure as a vector, not a scalar} --- secondary inorganics, EC, smoke
carbon and UFP number each carried a named effect in a large cohort; the leading constituent differs by
outcome (BC for stroke and heart failure in September, nitrate for lung cancer, sulfate/nitrate/ammonium for
kidney, EC for ADHD), which is either outcome specificity or model-to-model instability in the component
surfaces --- no record yet distinguishes the two. \textbf{(ii) Reference-side metrology} --- photometer
bias pooled with a prediction interval, filter-bound water, type-gated lidar mass, reference CPC drift,
housing-induced sample drying. \textbf{(iii) ML emulators and forecasters of the CTM} (ProNet across 338
cities; the Foehn generative km-scale forecaster; a CMAQ nitrate U-Net) --- the reason \textit{Exposure}
grew every week. \textbf{(iv) Respiratory endpoints} (1 $\rightarrow$ 17 per week), led by early-life lung
function (COPSAC FEV$_1$ $-$2.7\,\%, $-$4.27 to $-$1.1, at age 6), COPD and interstitial lung disease.
\textbf{(v) Non-exhaust traffic PM}: TRWP inventory, dynamometer tire-wear particles.

\vspace{2mm}
\textbf{Fading.} \textbf{Sensing share} 25 $\rightarrow$ 19 $\rightarrow$ \textbf{15\,\%} across August--October.
Part of this is the Crossref AMT/ACP channel and part is real: the instrumentation records that did appear
were reviews and validation audits rather than new devices. \textbf{Cardiovascular \& metabolic} halved
(46 $\rightarrow$ 28) and \textbf{tier A} fell from 23\,\% to \textbf{14\,\%} --- the latter partly a
metadata-only effect, since a record without an abstract cannot be tier A. Mass-only single-pollutant
time series continue to lose significance under co-pollutant adjustment; that observation has now held for
three months and is no longer news.

\vspace{2mm}
\textbf{Recurring groups and venues.} \textbf{No author-cluster analysis is possible}: the corpus still has
no affiliation field, and surname repeats are not investigator clusters. What does recur is a design
family --- national US claims cohorts (Medicare, Medicaid) with modelled smoke-fraction or component
exposure, which produced the ADRD, atrial-fibrillation and autism results. Venue concentration:
\textit{Environ Pollut} 27, \textit{Atmosphere} 24, \textit{Environ Sci Technol} 23, \textit{Atmos Environ}
23, \textit{Atmos Pollut Res} 20, \textit{Atmos Chem Phys} 20, \textit{Front Public Health} 15,
\textit{Build Environ} 13, \textit{Environ Res} 12, \textit{Environ Monit Assess} 12, \textit{Sensors} 8.
\textbf{Preprints rose to 33} (Research Square 22, medRxiv 11), so a larger share of the month is
unreviewed. \textbf{Most-cited work is not reported}: citation counts on records days to weeks old measure
indexing lag. The month's highest-venue records: Yu et al.\ (\textit{Lancet Planet Health}),
Papatheodorou et al.\ (\textit{JAMA Pediatr}), Manchanda et al.\ (\textit{Sci Adv}), and the two
\textit{Nat Commun} inventories (China BC burden; TRWP).

\clearpage

% ------------------------------------------------------------------ trial watch
\section*{\sffamily\bfseries\color{Deep}Trial registry watch}
\vspace{-1.5mm}{\color{Coral}\rule{\linewidth}{1.1pt}}\vspace{2.5mm}

\footnotesize
Registry sweeps (\texttt{LastUpdatePostDate} windows) ran through October with gaps covered by multi-day
windows; \texttt{state/trials.json} now holds \textbf{20 trials}, \textbf{four first seen in October}. The
pattern carried from August and September holds: \textbf{interventions framed around PM whose primary
endpoint is not a measured PM exposure.}

\vspace{1.5mm}
\textbf{NCT07536178} (``Lungs on Fire'', University of British Columbia, controlled exposure, $n$=24,
crossover; first seen 29 Oct): lodgepole-pine woodsmoke against filtered air in people with asthma, COPD or
lung cancer. \texttt{analyze\_endpoints}: co-primary \textbf{FeNO and CANTAB cognition}; secondaries
include sputum cells, resting-state fMRI default-mode connectivity, $^{129}$Xe MRI, a NULISAseq CNS panel,
NETs, epigenetic age and glial extracellular vesicles --- the first registered protocol in this watch
putting lung and brain endpoints under one smoke exposure, which is exactly the endpoint pairing the
October smoke epidemiology implies. \textbf{NCT02763917} (AIRWEIGHS, Johns Hopkins, $n$=198, placebo-
controlled home purifier in obese vs non-obese asthmatic children): record updated 21--22 Oct, primary
completion September 2026 --- \textbf{results posting is the event to watch}. \textbf{NCT07500948}
(EPIC-AIR, Leicester, feasibility RCT, $n$=80): real-time air-quality guidance during exercise
rehabilitation; the co-primary is \textbf{modelled} ambient PM at GPS location during activity, not a
personal measured dose, and there is no health co-primary. \textbf{NCT07886346} (Taif COPD prevalence
survey, observational): logged, low relevance.

\vspace{1.5mm}
Published trial-class records: \textbf{8 in 523}, the most relevant a randomised crossover of short-term
indoor PM$_{2.5}$ on cognition and HRV (\textit{iScience}, 26 Oct). The final sweep (30 Oct--1 Nov)
returned no hits. \textbf{Registry and literature now agree on one thing: the intervention evidence base
for PM is being built on filtration and smoke chambers, and the registered exposure endpoint, where one
exists, is increasingly modelled.}

\vspace{4mm}

% -------------------------------------------------------------------- open gaps
\section*{\sffamily\bfseries\color{Deep}Open gaps and read on direction}
\vspace{-1.5mm}{\color{Coral}\rule{\linewidth}{1.1pt}}\vspace{2.5mm}

\begin{keybox}
{\sffamily\bfseries\footnotesize\color{Deep}Gaps in the literature}\\[1.2mm]
\footnotesize
\textbf{(i) Component, number and smoke exposures are all modelled.} Not one October effect estimate on a
constituent, UFP number or smoke carbon rests on measured speciation, a CPC network or an absorption
channel. Model-to-model disagreement in component surfaces is unquantified, and it could produce the
outcome-specific ``leading constituent'' pattern on its own.
\textbf{(ii) Reference uncertainty is not propagated.} Calibrations report sensor error against a
reference treated as truth; October quantified photometer (+36\,\%), filter-water (+3 to +18\,\%) and
housing-RH biases, and no calibration paper carries them into its error budget.
\textbf{(iii) Number instruments still lack long records} --- no low-cost PN or CPC network with a
field calibration of one season or more appeared (September's gap (iii), unchanged).
\textbf{(iv) RH correction is fitted to the wrong humidity} for housed sensors; corrections developed on
bare PMS5003 units should not transfer to PurpleAir without re-fitting.
\textbf{(v) Trials register modelled or no exposure} (EPIC-AIR; three September registrations).
\end{keybox}

\vspace{2.5mm}

\begin{keybox}
{\sffamily\bfseries\footnotesize\color{Deep}Gaps and defects this pipeline creates}\\[1.2mm]
\footnotesize
\textbf{(i) Duplicates.} 523 records, \textbf{523 unique DOIs}. The 30 Oct backfill screened 23--29 Oct in
parallel and double-admitted $\sim$40 cross-day re-surfacings; the DOI gate and an earliest-issue rule
removed all of them before shipping, and backfills are now screened sequentially.
\textbf{(ii) 112 records (21\,\%) are metadata-only}, the third monthly rise (12, 16, 21\,\%). Scholar
Gateway returned access errors when called and was skipped on most runs --- it needs the user to reconnect it; Elsevier day-0 deposits
remain abstract-free.
\textbf{(iii) Run continuity.} Runs stalled for 6--9 Oct and 23--29 Oct; both gaps were backfilled day by
day with their own entry windows, so coverage is contiguous but several issues were built 1--7 days late.
\textbf{(iv) Source legs.} OpenAlex returned 429 on every run; Europe PMC showed day-0 index lag on several
days (0--3 records) and re-surfaces earlier PubMed records on later days; arXiv returned 150\,kB with no
in-window entry on most days; Consensus returned \textbf{0 new} records on every sweep.
\textbf{(v) Label maps.} Ecological designs split across two labels (f2a); the endpoint map is uncapped
(f2b omitted); no affiliation field, so no author clusters.
\end{keybox}

\vspace{3mm}

\footnotesize
\textbf{Scoring September's predictions.} (1) \textit{A low-cost network with an absorption/BC or
particle-number channel and a field calibration of at least one full season} --- \textbf{not met}.
The nearest record is Manchanda et al.'s West Oakland BC network (\textit{Sci Adv}, 28 Oct): network-scale
BC with Bayesian inversion, but not low-cost nodes and no season-long field calibration reported.
(2) \textit{A calibration paper reporting a temporally blocked or prospective hold-out as its primary
metric} --- \textbf{met}: Mercado-Velazquez et al.\ report seasonal hold-out (April--October training,
winter test) and cross-site transfer as headline results; the type-gated lidar used date-separated
validation, and a Nagoya PAH model used time-blocked CV.

\vspace{1.5mm}
\textbf{Read on direction.} August read the two halves of the field as optimising different objectives;
September narrowed the gap. October widened it again, from the health side: the epidemiology now names
components, number and source as the exposure, while the instruments that could measure those at
network scale are absent from the month's record and the instruments that measure mass were shown to be
calibrated against references with composition-dependent biases. The practical consequence for a
distributed PM network is that \textbf{density is no longer the binding constraint --- the measured
quantity and the reference are}.
\textbf{Falsifiable predictions for November:}
(1) a constituent-resolved or number-resolved health estimate whose exposure is \emph{measured}
(speciation network, CPC network, or network absorption/BC channel) rather than modelled; and
(2) a sensor-calibration paper that propagates reference-instrument uncertainty (gravimetric water,
FEM bias or photometer bias) into its reported error. If (2) fails, reported LCS errors should be read as
lower bounds.

\clearpage

% ==================================================================== DIGEST
%% MANDATORY, same contract as the daily and the weekly (formalised 2026-08-02).
%% Full summaries are lifted VERBATIM from issues/digest_*.tex by mkdigest.py.
\section*{\sffamily\bfseries\color{Deep}Paper-level digest --- the month by subtopic}
\vspace{-1.5mm}{\color{Coral}\rule{\linewidth}{1.1pt}}\vspace{2mm}

\footnotesize
Every pooled record appears below. Records that carry a full assessed summary have it lifted
verbatim from the daily issue that first published it; a title-only or extractive entry is
\textbf{not an assessment} and the absence of criticism is not approval. Metadata-only records are a
pointer to a DOI and nothing more.

\vspace{3mm}

\input{digest_body}

\clearpage

% ==================================================================== TABLE
\section*{\sffamily\bfseries\color{Deep}Pooled corpus register --- October 2026}
\vspace{-1.5mm}{\color{Coral}\rule{\linewidth}{1.1pt}}\vspace{2.5mm}

\input{table_rows}

\vspace{3mm}

\begin{keybox}
{\sffamily\bfseries\footnotesize\color{Deep}Provenance and method}\\[1.2mm]
{\fontsize{7.2}{8.8}\selectfont
Pure aggregation over \texttt{state/corpus/2026-10-*.json} (31 daily issues, entry windows 1--31
October 2026, contiguous), \texttt{state/metrics.csv} and \texttt{state/trials.json}. No API was
re-queried and no DOI re-resolved; every record passed the DOI gate at daily-issue time. Figures are
pooled re-renders of the daily set; the forest panels catalogue estimates and do not pool them.
Abstracts remain the copyright of their respective publishers.\par}
\end{keybox}

\end{document}
